The forced ``glide'' reentry (a known empirical acceleration injected on a
gravity-only baseline) is our EA test bed. Its exact construction is the
committed \code{glide\_fixture.json}, derived bit-for-bit from the canonical
internal scenario. With \emph{clean} 100\,Hz observations we recover the EA
several ways and re-propagate, RMSE vs truth.

Formally, every estimator below solves the regularized nonlinear
least-squares problem~\eqref{eq:opt} with the B-spline
parameterization~\eqref{eq:spline}; the knot spacing follows the spline
bias bound \eqref{eq:knot} (Theorem~\ref{thm:spline},
\cite{deboor2001}) and the GN refinement is the local
basin/continuation argument of Proposition~\ref{prop:basin}. See the
Mathematical Foundations chapter for the full theory.

\examplenote{\code{examples/cpp/ex04\_ea\_clean\_glide.cpp}}
\exfile{examples/cpp/ex04_ea_clean_glide.cpp}{C++}
\resultfile{results/ea_clean_glide.csv}{Clean-EA recovery by method.}

\textbf{Read every row against the first one.} The
\code{no\_ea\_baseline} row is the ``do nothing'' control: propagate the
glide IC with EA=0 and the trajectory is ${\sim}35\,$km off the truth ---
that is how much signal the empirical acceleration carries.

The lesson is then \textbf{method selection}. Bare B-spline POS
\emph{without} GN reprops at ${\sim}1.5\,$m --- already a ${\sim}24000\times$
improvement on the 35\,km baseline, and GN sharpens it another
${\sim}200\times$ to a few millimeters. GN is the load-bearing accuracy step
for position (it is on by default; this row switches it off to expose the
baseline). On clean 100\,Hz data velocity (${\sim}0.09\,$m) and acceleration
(${\sim}0.008\,$m) recover sub-meter too; the genuinely hard,
configuration-sensitive regime is \emph{noisy} data --- see Chapters 5 and 8.
